\documentclass[10pt,a4paper]{amsart}
\usepackage{amsmath,amssymb,amsthm,graphicx}
\usepackage{verbatim}
\newenvironment{reference}{\comment}{\endcomment}
%\newenvironment{reference}{}{}
\newenvironment{slogan}{\comment}{\endcomment}
\newenvironment{history}{\comment}{\endcomment}

% For commutative diagrams we use Xy-pic
\usepackage[all]{xy}

% We use 2cell for 2-commutative diagrams.
\xyoption{2cell}
% Theorem environments.
%
\theoremstyle{plain}
\newtheorem{theorem}[subsection]{Theorem}
\newtheorem{proposition}[subsection]{Proposition}
\newtheorem{lemma}[subsection]{Lemma}

\theoremstyle{definition}
\newtheorem{definition}[subsection]{Definition}
\newtheorem{example}[subsection]{Example}
\newtheorem{exercise}[subsection]{Exercise}
\newtheorem{situation}[subsection]{Situation}

\theoremstyle{remark}
\newtheorem{remark}[subsection]{Remark}
\newtheorem{remarks}[subsection]{Remarks}

\numberwithin{equation}{subsection}

% Macros
%
\def\lim{\mathop{\mathrm{lim}}\nolimits}
\def\colim{\mathop{\mathrm{colim}}\nolimits}
\def\Spec{\mathop{\mathrm{Spec}}}
\def\Hom{\mathop{\mathrm{Hom}}\nolimits}
\def\Ext{\mathop{\mathrm{Ext}}\nolimits}
\def\SheafHom{\mathop{\mathcal{H}\!\mathit{om}}\nolimits}
\def\SheafExt{\mathop{\mathcal{E}\!\mathit{xt}}\nolimits}
\def\Sch{\mathit{Sch}}
\def\Mor{\mathop{\mathrm{Mor}}\nolimits}
\def\Ob{\mathop{\mathrm{Ob}}\nolimits}
\def\Sh{\mathop{\mathit{Sh}}\nolimits}
\def\NL{\mathop{N\!L}\nolimits}
\def\CH{\mathop{\mathrm{CH}}\nolimits}
\def\proetale{{pro\text{-}\acute{e}tale}}
\def\etale{{\acute{e}tale}}
\def\QCoh{\mathit{QCoh}}
\def\Ker{\mathop{\mathrm{Ker}}}
\def\Im{\mathop{\mathrm{Im}}}
\def\Coker{\mathop{\mathrm{Coker}}}
\def\Coim{\mathop{\mathrm{Coim}}}

% Boxtimes
%
\DeclareMathSymbol{\boxtimes}{\mathbin}{AMSa}{"02}

%
% Macros for moduli stacks/spaces
%
\def\QCohstack{\mathcal{QC}\!\mathit{oh}}
\def\Cohstack{\mathcal{C}\!\mathit{oh}}
\def\Spacesstack{\mathcal{S}\!\mathit{paces}}
\def\Quotfunctor{\mathrm{Quot}}
\def\Hilbfunctor{\mathrm{Hilb}}
\def\Curvesstack{\mathcal{C}\!\mathit{urves}}
\def\Polarizedstack{\mathcal{P}\!\mathit{olarized}}
\def\Complexesstack{\mathcal{C}\!\mathit{omplexes}}
% \Pic is the operator that assigns to X its picard group, usage \Pic(X)
% \Picardstack_{X/B} denotes the Picard stack of X over B
% \Picardfunctor_{X/B} denotes the Picard functor of X over B
\def\Pic{\mathop{\mathrm{Pic}}\nolimits}
\def\Picardstack{\mathcal{P}\!\mathit{ic}}
\def\Picardfunctor{\mathrm{Pic}}
\def\Deformationcategory{\mathcal{D}\!\mathit{ef}}
\usepackage{cleveref}
\crefname{originallemma}{original lemma}{original lemmas}
\Crefname{originallemma}{Original lemma}{Original lemmas}
\makeatletter
\newlabel{algebra-lemma-completion-generalities}{{0315}{source}{}{sourcepointer-1}{}}
\newlabel{algebra-lemma-completion-generalities@cref}{{[originallemma][0][]0315}{[1][0][]source}{}{}{}}
\makeatother
\crefname{originallemma}{original lemma}{original lemmas}
\Crefname{originallemma}{Original lemma}{Original lemmas}
\makeatletter
\newlabel{algebra-lemma-intersect-powers-ideal-module-zero}{{00IP}{source}{}{sourcepointer-2}{}}
\newlabel{algebra-lemma-intersect-powers-ideal-module-zero@cref}{{[originallemma][0][]00IP}{[1][0][]source}{}{}{}}
\makeatother
\crefname{originallemma}{original lemma}{original lemmas}
\Crefname{originallemma}{Original lemma}{Original lemmas}
\makeatletter
\newlabel{algebra-lemma-surjective-endo-noetherian-ring-is-iso}{{06RN}{source}{}{sourcepointer-3}{}}
\newlabel{algebra-lemma-surjective-endo-noetherian-ring-is-iso@cref}{{[originallemma][0][]06RN}{[1][0][]source}{}{}{}}
\makeatother
\crefname{originalreference}{original reference}{original references}
\Crefname{originalreference}{Original reference}{Original references}
\makeatletter
\newlabel{definition-rig-smooth-homomorphism}{{0GAI}{source}{}{sourcepointer-4}{}}
\newlabel{definition-rig-smooth-homomorphism@cref}{{[originalreference][0][]0GAI}{[1][0][]source}{}{}{}}
\makeatother
\crefname{originalreference}{original reference}{original references}
\Crefname{originalreference}{Original reference}{Original references}
\makeatletter
\newlabel{lemma-get-morphism-general-better}{{0GAQ}{source}{}{sourcepointer-5}{}}
\newlabel{lemma-get-morphism-general-better@cref}{{[originalreference][0][]0GAQ}{[1][0][]source}{}{}{}}
\makeatother
\crefname{originalsection}{original section}{original sections}
\Crefname{originalsection}{Original section}{Original sections}
\makeatletter
\newlabel{local-cohomology-section-uniform}{{0GA6}{source}{}{sourcepointer-6}{}}
\newlabel{local-cohomology-section-uniform@cref}{{[originalsection][0][]0GA6}{[1][0][]source}{}{}{}}
\makeatother
\crefname{originallemma}{original lemma}{original lemmas}
\Crefname{originallemma}{Original lemma}{Original lemmas}
\makeatletter
\newlabel{more-algebra-lemma-approximate-complex}{{07VE}{source}{}{sourcepointer-7}{}}
\newlabel{more-algebra-lemma-approximate-complex@cref}{{[originallemma][0][]07VE}{[1][0][]source}{}{}{}}
\makeatother
\crefname{originallemma}{original lemma}{original lemmas}
\Crefname{originallemma}{Original lemma}{Original lemmas}
\makeatletter
\newlabel{more-algebra-lemma-approximate-complex-graded}{{07VF}{source}{}{sourcepointer-8}{}}
\newlabel{more-algebra-lemma-approximate-complex-graded@cref}{{[originallemma][0][]07VF}{[1][0][]source}{}{}{}}
\makeatother
\title{Rig-smooth quotient algebras are stable under sufficiently close deformations of their presentations}
\author{The Stacks Project Authors}
\date{stacks-project:restricted.tex:local-snapshot-20261012; GNU FDL 1.2 or later}
\begin{document}
\maketitle
\noindent\textit{Editorial note.} All delivered mathematical text is the author's own, copied from the pinned Stacks Project chapter file restricted.tex; no statement or proof step was rewritten, reordered or removed.\par\medskip
\noindent\textit{Editorial note.} The counted proof is the author's complete argument: a three-constant bound n, a lifting of the mod-I\textasciicircum{}n comparison to a genuine homomorphism, and an isomorphism proof combining surjectivity by completeness with injectivity via the associated graded and Krull intersection, all inside the delivered span\par\medskip
\noindent\textit{Editorial note.} The label lemma-close-enough names the counted statement, and the proof environment immediately following it in the source is delivered verbatim as its proof.\par\medskip
\noindent\textit{Editorial note.} The counted statement is carried by the author's own lemma environment in the chapter source.\par\medskip
\noindent\textit{Editorial note.} Before the final part of the argument the author writes 'we suggest the reader find their own proof of this' and then immediately provides the complete proof (surjectivity, associated graded, Krull intersection); per the quality criterion this reading suggestion is not an omission, and the verification is delivered verbatim.\par\medskip
\noindent\textit{Allowed standard prerequisites (editorial).}\begin{itemize}
\item The rig-smooth homomorphism definition via annihilation of Ext\textasciicircum{}1 of the completed naive cotangent complex
\item The approximation lemmas for complexes over adic rings from More on Algebra, giving exactness of the perturbed presentation and an isomorphism of associated graded rings for n large
\item The uniformity integers d and q of the associated graded, and the lifting lemma producing a homomorphism congruent to a prescribed one modulo a prescribed power
\end{itemize}\medskip
\ifdefined\hypertarget\hypertarget{sourcepointer-1}{}\fi
\noindent\textit{Source reference.} Original lemma 0315: sources/stacks-chow/tags.txt, source lines 3972--3972. Prerequisite pointer invoked at proof line(s) 1434; it names a standard fact of a sibling chapter that the counted argument uses, and every step of the counted proof built on it is retained verbatim in the delivered body.\par\medskip
\ifdefined\hypertarget\hypertarget{sourcepointer-2}{}\fi
\noindent\textit{Source reference.} Original lemma 00IP: sources/stacks-chow/tags.txt, source lines 687--687. Prerequisite pointer invoked at proof line(s) 1447; it names a standard fact of a sibling chapter that the counted argument uses, and every step of the counted proof built on it is retained verbatim in the delivered body.\par\medskip
\ifdefined\hypertarget\hypertarget{sourcepointer-3}{}\fi
\noindent\textit{Source reference.} Original lemma 06RN: sources/stacks-chow/tags.txt, source lines 8647--8647. Prerequisite pointer invoked at proof line(s) 1439; it names a standard fact of a sibling chapter that the counted argument uses, and every step of the counted proof built on it is retained verbatim in the delivered body.\par\medskip
\ifdefined\hypertarget\hypertarget{sourcepointer-4}{}\fi
\noindent\textit{Source reference.} Original reference 0GAI: sources/stacks-restricted/pinned-original-source.tex, source lines 831--831. Prerequisite pointer invoked at proof line(s) 1383; it names a standard fact of the same chapter that the counted argument uses, and every step of the counted proof built on it is retained verbatim in the delivered body.\par\medskip
\ifdefined\hypertarget\hypertarget{sourcepointer-5}{}\fi
\noindent\textit{Source reference.} Original reference 0GAQ: sources/stacks-restricted/pinned-original-source.tex, source lines 1167--1167. Prerequisite pointer invoked at proof line(s) 1419; it names a standard fact of the same chapter that the counted argument uses, and every step of the counted proof built on it is retained verbatim in the delivered body.\par\medskip
\ifdefined\hypertarget\hypertarget{sourcepointer-6}{}\fi
\noindent\textit{Source reference.} Original section 0GA6: sources/stacks-chow/tags.txt, source lines 20326--20326. Prerequisite pointer invoked at proof line(s) 1401; it names a standard fact of a sibling chapter that the counted argument uses, and every step of the counted proof built on it is retained verbatim in the delivered body.\par\medskip
\ifdefined\hypertarget\hypertarget{sourcepointer-7}{}\fi
\noindent\textit{Source reference.} Original lemma 07VE: sources/stacks-chow/tags.txt, source lines 10004--10004. Prerequisite pointer invoked at proof line(s) 1388; it names a standard fact of a sibling chapter that the counted argument uses, and every step of the counted proof built on it is retained verbatim in the delivered body.\par\medskip
\ifdefined\hypertarget\hypertarget{sourcepointer-8}{}\fi
\noindent\textit{Source reference.} Original lemma 07VF: sources/stacks-chow/tags.txt, source lines 10005--10005. Prerequisite pointer invoked at proof line(s) 1389; it names a standard fact of a sibling chapter that the counted argument uses, and every step of the counted proof built on it is retained verbatim in the delivered body.\par\medskip
\begin{lemma}
\label{lemma-close-enough}
Let $I$ be an ideal of a Noetherian ring $A$. Let $r \geq 0$
and write $P = A[x_1, \ldots, x_r]$ the $I$-adic completion.
Consider a resolution
$$
P^{\oplus t} \xrightarrow{K} P^{\oplus m}
\xrightarrow{g_1, \ldots, g_m} P \to B \to 0
$$
of a quotient of $P$. Assume $B$ is rig-smooth over $(A, I)$.
Then there exists an integer $n$ such that for any complex
$$
P^{\oplus t} \xrightarrow{K'} P^{\oplus m}
\xrightarrow{g'_1, \ldots, g'_m} P
$$
with $g_i - g'_i \in I^nP$ and $K - K' \in I^n\text{Mat}(m \times t, P)$
there exists an isomorphism $B \to B'$ of $A$-algebras where
$B' = P/(g'_1, \ldots, g'_m)$.
\end{lemma}

\begin{proof}
(A) By Definition \ref{definition-rig-smooth-homomorphism}
we can choose a $c \geq 0$ such that $I^c$ annihilates
$\Ext^1_B(\NL_{B/A}^\wedge, N)$ for all $B$-modules $N$.

\medskip\noindent
(B) By More on Algebra, Lemmas \ref{more-algebra-lemma-approximate-complex} and
\ref{more-algebra-lemma-approximate-complex-graded}
there exists a constant $c_1 = c(g_1, \ldots, g_m, K)$
such that for $n \geq c_1 + 1$ the complex
$$
P^{\oplus t} \xrightarrow{K'} P^{\oplus m}
\xrightarrow{g'_1, \ldots, g'_m} P \to B' \to 0
$$
is exact and $\text{Gr}_I(B) \cong \text{Gr}_I(B')$.

\medskip\noindent
(C) Let $d_0 = d(\text{Gr}_I(B))$ and $q_0 = q(\text{Gr}_I(B))$
be the integers found in
Local Cohomology, Section \ref{local-cohomology-section-uniform}.

\medskip\noindent
We claim that $n = \max(c_1 + 1, q_0 + (d_0 + 1)c, 2(d_0 + 1)c + 1)$
works where $c$ is as in (A), $c_1$ is as in (B), and $q_0, d_0$ are as
in (C).

\medskip\noindent
Let $g'_1, \ldots, g'_m$ and $K'$ be as in the lemma.
Since $g_i = g'_i \in I^nP$ we obtain a canonical
$A$-algebra homomorphism
$$
\psi_n : B \longrightarrow B'/I^nB'
$$
which induces an isomorphism $B/I^nB \to B'/I^nB'$. Since
$\text{Gr}_I(B) \cong \text{Gr}_I(B')$ we have
$d_0 = d(\text{Gr}_I(B'))$ and $q_0 = q(\text{Gr}_I(B'))$
and since $n \geq \max(q_0 + (1 + d_0)c, 2(d_0 + 1)c + 1)$
we may apply Lemma \ref{lemma-get-morphism-general-better}
to find an $A$-algebra homomorphism
$$
\varphi : B \longrightarrow B'
$$
such that
$\varphi \bmod I^{n - (d_0 + 1)c}B' = \psi_n \bmod I^{n - (d_0 + 1)c}B'$.
Since $n - (d_0 + 1)c > 0$ we see that $\varphi$ is an $A$-algebra
homomorphism which modulo $I$ induces the isomorphism $B/IB \to B'/IB'$
we found above.
The rest of the proof shows that these facts force $\varphi$
to be an isomorphism; we suggest the reader find their own proof of this.

\medskip\noindent
Namely, it follows that $\varphi$ is surjective for example by applying
Algebra, Lemma \ref{algebra-lemma-completion-generalities} part (1)
using the fact that $B$ and $B'$ are complete.
Thus $\varphi$ induces a surjection $\text{Gr}_I(B) \to \text{Gr}_I(B')$
which has to be an isomorphism because the source and target are
isomorphic Noetherian rings, see
Algebra, Lemma \ref{algebra-lemma-surjective-endo-noetherian-ring-is-iso}
(of course you can show $\varphi$ induces the isomorphism we found above
but that would need a tiny argument).
Thus $\varphi$ induces injective maps
$I^eB/I^{e + 1}B \to I^eB'/I^{e + 1}B'$ for all $e \geq 0$.
This implies $\varphi$ is injective since for any $b \in B$ there
exists an $e \geq 0$ such that $b \in I^eB$, $b \not \in I^{e + 1}B$
by Krull's intersection theorem (Algebra, Lemma
\ref{algebra-lemma-intersect-powers-ideal-module-zero}).
This finishes the proof.
\end{proof}

\end{document}
